% SPDX-FileCopyrightText: 2026 Bart van de Lint
% SPDX-License-Identifier: AGPL-3.0-or-later
\section{Ship classes}
\label{sec:classes}

\Cref{tab:classes} gives the thirteen ship classes of \cref{sec:ships}. The source of each value is in \cref{tab:class_sources} and every other input, with its range and source, in \cref{tab:inputs}.

\begin{table*}[tbp]
\centering
\footnotesize
\caption{\stale{sensitivity}{} The thirteen standard ship classes, ordered by deadweight as in all class figures: length between perpendiculars $L$, beam $B$, design draught $T$, block coefficient $C_B$ at that draught (the underwater volume over that of a box of the same length, beam and draught), deadweight at the scantling draught, slot or deck capacity of the space-limited ships in TEU or standard cars (CEU), the limit the cargo of the class fills first, service speed and main-engine power of the source design, the factor $f_R$ that calibrates the calm-water resistance to it, newbuild price, cargo value, container tiers on deck, the lowest air draft (height above the waterline) on common calls, the share of calls a fixed rig above it loses and the time a fixed rig costs cargo handling per call ($\infty$ where it prevents it), and the frontal and lateral areas above the waterline with the deck cargo, our estimates scaled to each class from the measured ship of its type \citep{blendermann1993} and given to three digits; \cref{sec:deck_rigs} explains the air draft, the lost calls and the delay. Where a source gives only the scantling draught, as for the container ships and the Aframax tanker, the design draught is the scantling draught times the ratio of the two in a published design of the same type. Each value has its source in \cref{tab:class_sources}; those marked $\dagger$ are our estimates, for which no range is given; the coaster's cargo value is the cargo value input of \cref{tab:inputs}, given with its range.}
\label{tab:classes}
\setlength{\tabcolsep}{3pt}
\begin{tabular}{@{}lrrrrrrlrrrr@{}}
\toprule
Class & $L$ & $B$ & $T$ & $C_B$ & Deadweight & Capacity & Limit & Speed & Power & $f_R$ & Price \\
 & [m] & [m] & [m] & & [t] & & & [kn] & [MW] & & [MUSD] \\
\midrule
\multicolumn{12}{@{}l}{\textit{Free deck}} \\
general cargo coaster & \num{85} & \num{13.2} & \num{5.5} & \num{0.83} & \num{3800} & -- & deadweight & \num{10.5} & \num{1.0} & \num{1.00} & \num{11.6} \\
PCTC car carrier & \num{192} & \num{37.6} & \num{9.3} & \num{0.62} & \num{25563} & \num{9100} CEU & space & \num{14.5} & \num{13.7} & \num{0.97} & \num{100.0} \\
Handysize bulk carrier & \num{177} & \num{30.5} & \num{10.3} & \num{0.80} & \num{42000} & -- & deadweight & \num{14.0} & \num{6.5} & \num{0.93} & \num{31.0} \\
MR product tanker & \num{174} & \num{32.2} & \num{11.0} & \num{0.81} & \num{47000} & -- & deadweight & \num{15.1} & \num{9.3} & \num{0.96} & \num{51.0} \\
Ultramax bulk carrier & \num{196} & \num{32.2} & \num{12.2} & \num{0.86} & \num{64000} & -- & deadweight & \num{14.5} & \num{8.6} & \num{0.80} & \num{34.5} \\
Kamsarmax bulk carrier & \num{225} & \num{32.3} & \num{12.2} & \num{0.88} & \num{82000} & -- & deadweight & \num{14.5} & \num{9.7} & \num{0.82} & \num{37.5} \\
Aframax tanker & \num{236} & \num{44.0} & \num{13.9} & \num{0.80} & \num{114824} & -- & deadweight & \num{14.2} & \num{10.3} & \num{0.92} & \num{75.0} \\
Capesize bulk carrier & \num{282} & \num{45.0} & \num{16.5} & \num{0.84} & \num{180000} & -- & deadweight & \num{14.6} & \num{16.2} & \num{0.87} & \num{75.0} \\
VLCC & \num{320} & \num{60.0} & \num{20.5} & \num{0.81} & \num{301781} & -- & deadweight & \num{15.5} & \num{27.2} & \num{1.05} & \num{130.0} \\
\multicolumn{12}{@{}l}{\textit{Cargo on deck}} \\
feeder container ship & \num{165} & \num{25.6} & \num{8.8} & \num{0.66} & \num{21000} & \num{1600} TEU & space & \num{18.9} & \num{10.5} & \num{0.93} & \num{33.0} \\
Panamax container ship & \num{245} & \num{32.2} & \num{11.2} & \num{0.65} & \num{51000} & \num{4000} TEU & space & \num{21.1} & \num{22.8} & \num{0.96} & \num{58.0} \\
Neo-Panamax container ship & \num{350} & \num{50.8} & \num{14.3} & \num{0.67} & \num{152000} & \num{14000} TEU & space & \num{22.4} & \num{49.4} & \num{0.93} & \num{152.5} \\
Megamax container ship & \num{398} & \num{61.5} & \num{14.8} & \num{0.72} & \num{230000} & \num{24000} TEU & deadweight & \num{21.7} & \num{60.0} & \num{0.90} & \num{267.3} \\
\bottomrule
\end{tabular}

\medskip
\begin{tabular}{@{}lrrrrrrr@{}}
\toprule
Class & Cargo value & Deck & Air draft$^\dagger$ & Lost$^\dagger$ & Fixed-rig$^\dagger$ & $A_\mathrm{F}$ & $A_\mathrm{L}$ \\
 & [USD/t] & tiers & limit [m] & calls [\%] & delay [h] & [m$^2$] & [m$^2$] \\
\midrule
\multicolumn{8}{@{}l}{\textit{Free deck}} \\
general cargo coaster & \num{2500} (\numrange{1900}{3700}) & \num{0} & \num{35} & \num{10} & $\infty$ & \num{164} & \num{598} \\
PCTC car carrier & \num{15000} & \num{0} & \num{53} & \num{5} & \num{0} & \num{1090} & \num{5220} \\
Handysize bulk carrier & \num{250}\rlap{$^\dagger$} & \num{0} & \num{45} & \num{6} & \num{2} & \num{328} & \num{944} \\
MR product tanker & \num{650}\rlap{$^\dagger$} & \num{0} & \num{45} & \num{6} & \num{0} & \num{341} & \num{912} \\
Ultramax bulk carrier & \num{180}\rlap{$^\dagger$} & \num{0} & \num{45} & \num{6} & \num{2} & \num{384} & \num{1160} \\
Kamsarmax bulk carrier & \num{170}\rlap{$^\dagger$} & \num{0} & \num{45} & \num{6} & \num{4} & \num{441} & \num{1530} \\
Aframax tanker & \num{506} & \num{0} & \num{45} & \num{6} & \num{0} & \num{631} & \num{1680} \\
Capesize bulk carrier & \num{100} & \num{0} & \num{60} & \num{1} & \num{4} & \num{772} & \num{2400} \\
VLCC & \num{506} & \num{0} & \num{62} & \num{1} & \num{0} & \num{1170} & \num{3090} \\
\multicolumn{8}{@{}l}{\textit{Cargo on deck}} \\
feeder container ship & \num{5340} & \num{4} & \num{40} & \num{5} & \num{0} & \num{623} & \num{2950} \\
Panamax container ship & \num{5340} & \num{6} & \num{53} & \num{5} & \num{0} & \num{1170} & \num{6540} \\
Neo-Panamax container ship & \num{5340} & \num{8} & \num{53} & \num{5} & \num{0} & \num{2560} & \num{13000} \\
Megamax container ship & \num{5340} & \num{10} & \num{57} & \num{5} & \num{0} & \num{3660} & \num{17500} \\
\bottomrule
\end{tabular}

\end{table*}
